\begin{enumerate}

\item  We have
  \begin{align*}
\lim_{h\to 0} \frac{\frac{x+h}{x+h-1}-\frac{x}{x-1}}{h}&= \lim_{h\to
  0} \frac{\frac{x^2-x+hx-h-x^2-hx+x}{(x+h-1)(x-1)}}{h} \\&=
\lim_{h\to 0} \frac{-1}{(x+h-1)(x-1)} \\&= \frac{-1}{(x-1)^2}.	
  \end{align*}
The limit exists for all $x$ except $x=1$.


\item If $f: \mathbb{R}^n \rightarrow \mathbb{R}$ is a differentiable
function, the function $\nabla f(x): \mathbb{R}^n \rightarrow \mathbb{R}^n$ is called the gradient of $f$ and is defined as 
\[
  \nabla f(x_1,x_2,x_3)
=
\cvect{
     \frac{\partial f}{\partial x_1} \\
    \frac{\partial f}{\partial x_2}  \\
    \frac{\partial f}{\partial x_3} 
}.
\]
Therefore, the gradient of our function is
\[
 \nabla f(x_1,x_2,x_3)
=
\cvect{
     e^{x_1}+2x_1x_3-x_2x_3 \\
    -x_1x_3 \\
    x_1^2-x_1x_2
}.
\]

\item  Let $f: \mathbb{R}^n \rightarrow \mathbb{R}$ be a differentiable
 function. Consider $x \in \mathbb{R}^n$ and $d \in \mathbb{R}^n$. The
 directional derivative of $f$ in $x$ in the direction $d$ is given by
 \[
\lim_{\alpha\to 0} \frac{f(x+ \alpha d)-f(x)}{\alpha}, 
\]
if the limit exists. In addition, when the gradient exists, the
directional derivative is the scalar product between the gradient of
$f$ and the direction $d$, that is,
\[
\nabla f(x)^Td.
\]
Therefore, the directional derivative is
\begin{align*}
(d_1 \, d_2 \,d_3)\nabla f(x_1,x_2,x_3)=
\end{align*}
\[d_1(e^{x_1}+2x_1x_3-x_2x_3)-d_2x_1x_3-d_3(x_1^2-x_1x_2). \]

\item This is incorrect. Indeed, the gradient is
$$\nabla f(x_1,x_2)= 
\begin{pmatrix}
2x_1\\
2x_2
\end{pmatrix}
$$ 
and the point $(0,0)$, we obtain as directional derivative:
$$\text{$\nabla$}f(0,0)^Td=(0,0)\begin{pmatrix}
1\\1
\end{pmatrix}=0.$$


\item Recall that the direction opposite the gradient is that in which the function has its steepest descent.
$$\nabla f(x)
=
\begin{pmatrix}
\frac{\partial f}{\partial x_1}\\
\frac{\partial f}{\partial x_2}
\end{pmatrix}
=
\begin{pmatrix}
x_1\\
4x_2
\end{pmatrix}.$$

In point $x=(1 \, 1)^T$
\[
\nabla f(x)
=
\begin{pmatrix}
x_1\\
4x_2
\end{pmatrix}
=
\begin{pmatrix}
1\\
4
\end{pmatrix},
\]

The directional derivative in f in the direction opposite to the gradient is:

\begin{align*}
-\nabla f(x)^T \nabla f(x)= \begin{pmatrix}
-1 & -4
\end{pmatrix}
\begin{pmatrix}
1\\
4
\end{pmatrix}
=
-17\\
\end{align*}

\item First, we are going to calculate the gradient matrix:
\[
\nabla f(x)
=
\begin{pmatrix}
   \frac{\partial f_1}{\partial x_1}      & \frac{\partial f_2}{\partial x_1} \\
    \frac{\partial f_1}{\partial x_2}       & \frac{\partial f_2}{\partial x_2}
\end{pmatrix}
=
\begin{pmatrix}
   2x_1x_2      & -\sin(e^{x_1+x_2^2})e^{x_1+x_2^2} \\
    x_1^2      & -2x_2\sin(e^{x_1+x_2^2})e^{x_1+x_2^2}
    \end{pmatrix}.
\]

Knowing that transposing the gradient we obtain the Jacobian, we are going to calculate it
\[
   JF(x_1,x_2)=\nabla f(x)^T
=
\begin{pmatrix}
   \frac{\partial f_1}{\partial x_1}      & \frac{\partial f_1}{\partial x_2} \\
    \frac{\partial f_2}{\partial x_1}       & \frac{\partial f_2}{\partial x_2}
\end{pmatrix}
=
\begin{pmatrix}
   2x_1x_2      &   x_1^2 \\
    -\sin(e^{x_1+x_2^2})e^{x_1+x_2^2}    & -2x_2\sin(e^{x_1+x_2^2})e^{x_1+x_2^2}
\end{pmatrix}.
\]



\item The function $f(x,y)=400-3x^2y^2$ is differentiable in point $(1,-2)$, therefore the gradient is going to give us the direction in which the depth has its steepest ascent.
$$\text{$\nabla$}f(x,y)=(-6xy^2,-6x^2y)$$
$$\text{$\nabla$}f(1,-2)=(-24,12).$$
The normalized direction in which the depth increases as quick as possible is then $(-2/\sqrt{5},1/\sqrt{5})$.
\end{enumerate}

